%% LyX 2.4.4 created this file.  For more info, see https://www.lyx.org/.
%% Do not edit unless you really know what you are doing.
\documentclass[oneside,english]{amsart}
\usepackage[T1]{fontenc}
\usepackage[latin9]{inputenc}
\synctex=-1
\usepackage{mathrsfs}
\usepackage{amsthm}

\makeatletter
%%%%%%%%%%%%%%%%%%%%%%%%%%%%%% Textclass specific LaTeX commands.
\newlength{\lyxlabelwidth}      % auxiliary length 
\theoremstyle{definition}
\newtheorem*{defn*}{\protect\definitionname}
\theoremstyle{plain}
\newtheorem{thm}{\protect\theoremname}
\theoremstyle{plain}
\newtheorem{prop}[thm]{\protect\propositionname}

%%%%%%%%%%%%%%%%%%%%%%%%%%%%%% User specified LaTeX commands.
\date{}
\usepackage{commath}

\makeatother

\usepackage{babel}
\providecommand{\definitionname}{Definition}
\providecommand{\propositionname}{Proposition}
\providecommand{\theoremname}{Theorem}

\begin{document}
\begin{defn*}
A \emph{cell product function} is a real valued function whose domain
is the set of all binary lattices and whose value on any lattice is
the product of its values on every cell in the lattice.
\end{defn*}
{[}{[}I'm proposing what's below as a theorem rather than proposition,
to be parallel to Theorem~\ref{thm: sum of vs}. I feel both should
be theorems or both propositions, and Theorem~\ref{thm: sum of vs}
seems worthy of being a theorem rather than a proposition.{]}{]}
\begin{thm}
Let $v$ be the cell product function that maps the cells $\begin{smallmatrix}0 & 1\\
1 & 0
\end{smallmatrix}$ and $\begin{smallmatrix}1 & 0\\
0 & 1
\end{smallmatrix}$ to $0$, and maps the fourteen other cells cells to $1$. Then for
all positive integers $m$ and $n$, the number of $\left(m,n\right)$
polygon mosaics equals $\sum_{\ell}v\left(\ell\right)$, where the
sum is over all $\left(m,n\right)$ framed lattices $\ell$.
\end{thm}

\begin{proof}
For any framed lattice $\ell$, $v\left(\ell\right)$ is $1$ if $\ell$
is proper, i.e., has neither $\begin{smallmatrix}0 & 1\\
1 & 0
\end{smallmatrix}$ nor $\begin{smallmatrix}1 & 0\\
0 & 1
\end{smallmatrix}$cells, and otherwise is $0$, so that $\sum_{\ell}v\left(\ell\right)$
counts the number of proper $\left(m,n\right)$ framed lattices. Our
conclusion now follows from the fact that $\mathscr{L}$ gives a bijection
from the set of $\left(m,n\right)$ polygon mosaics to the set of
proper $\left(m,n\right)$ framed lattices.
\end{proof}
{[}{[}Here is a proposed rephrasing of material that on April~8 was
in Section~4. For a major result like this, I think it might be better
to put everything in one statement rather than spreading it through
prior definitions.{]}{]}
\begin{thm}
Let $v$ be any cell product function. Define the $1\times1$ matrices
$A_{1}=\begin{pmatrix}v\left(\begin{smallmatrix}0 & 0\\
0 & 0
\end{smallmatrix}\right)\end{pmatrix}$, $B_{1}=\begin{pmatrix}v\left(\begin{smallmatrix}0 & 0\\
1 & 0
\end{smallmatrix}\right)\end{pmatrix}$, $C_{1}=\begin{pmatrix}v\left(\begin{smallmatrix}1 & 0\\
0 & 0
\end{smallmatrix}\right)\end{pmatrix}$, and $D_{1}=\begin{pmatrix}v\left(\begin{smallmatrix}1 & 0\\
1 & 0
\end{smallmatrix}\right)\end{pmatrix}$, and for positive integers $n$, define the $2^{n}\times2^{n}$ matrices
{[}{[}We used to say $2^{n-1}\times2^{n-1}$; this bothered me for
a while, and now I realize it was wrong!! The issue is that while
$A_{n}$, $B_{n}$, etc. are indeed $2^{n-1}\times2^{n-1}$, here
we are defining $A_{n+1}$, $B_{n+1}$, etc.{]}{]}
\begin{eqnarray*}
A_{n+1}=\begin{pmatrix}v\left(\begin{smallmatrix}0 & 0\\
0 & 0
\end{smallmatrix}\right)A_{n} & v\left(\begin{smallmatrix}0 & 0\\
0 & 1
\end{smallmatrix}\right)B_{n}\\
v\left(\begin{smallmatrix}0 & 1\\
0 & 0
\end{smallmatrix}\right)C_{n} & v\left(\begin{smallmatrix}0 & 1\\
0 & 1
\end{smallmatrix}\right)D_{n}
\end{pmatrix} & B_{n+1}=\begin{pmatrix}v\left(\begin{smallmatrix}0 & 0\\
1 & 0
\end{smallmatrix}\right)A_{n} & v\left(\begin{smallmatrix}0 & 0\\
1 & 1
\end{smallmatrix}\right)B_{n}\\
v\left(\begin{smallmatrix}0 & 1\\
1 & 0
\end{smallmatrix}\right)C_{n} & v\left(\begin{smallmatrix}0 & 1\\
1 & 1
\end{smallmatrix}\right)D_{n}
\end{pmatrix}\\
C_{n+1}=\begin{pmatrix}v\left(\begin{smallmatrix}1 & 0\\
0 & 0
\end{smallmatrix}\right)A_{n} & v\left(\begin{smallmatrix}1 & 0\\
0 & 1
\end{smallmatrix}\right)B_{n}\\
v\left(\begin{smallmatrix}1 & 1\\
0 & 0
\end{smallmatrix}\right)C_{n} & v\left(\begin{smallmatrix}1 & 1\\
0 & 1
\end{smallmatrix}\right)D_{n}
\end{pmatrix} & D_{n+1}=\begin{pmatrix}v\left(\begin{smallmatrix}1 & 0\\
1 & 0
\end{smallmatrix}\right)A_{n} & v\left(\begin{smallmatrix}1 & 0\\
1 & 1
\end{smallmatrix}\right)B_{n}\\
v\left(\begin{smallmatrix}1 & 1\\
1 & 0
\end{smallmatrix}\right)C_{n} & v\left(\begin{smallmatrix}1 & 1\\
1 & 1
\end{smallmatrix}\right)D_{n}
\end{pmatrix}.
\end{eqnarray*}
 Then for all positive integers $m$ and $n$, the upper left entry
of $A_{n}^{m}$ equals $\sum_{\ell}v\left(\ell\right)$, where the
sum is over all $\left(m,n\right)$ framed lattices $\ell$. \label{thm: m by n matrix}
{[}{[}I changed ``$\left(0,0\right)$ entry'' to ``upper left entry''
because it means a casual reader doesn't need to know how we're numbering
matrix entries to understand this theorem's conclusion. Thoughts?{]}{]}
\end{thm}


\section{Counting polygon-free Mosaics}

\begin{thm}
Let $v$ be the cell product function that maps the cells $\begin{smallmatrix}0 & 1\\
1 & 0
\end{smallmatrix}$ and $\begin{smallmatrix}1 & 0\\
0 & 1
\end{smallmatrix}$ to $0$, $\begin{smallmatrix}0 & 0\\
0 & 0
\end{smallmatrix}$ and $\begin{smallmatrix}1 & 1\\
1 & 1
\end{smallmatrix}$ to $7$, $\begin{smallmatrix}0 & 0\\
1 & 0
\end{smallmatrix}$ and $\begin{smallmatrix}1 & 0\\
1 & 1
\end{smallmatrix}$ to $-1$, and the ten other cells to $1$. Then for all positive
integers $m$ and $n$, the number of $\left(m,n\right)$ polygon-free
mosaics equals $\sum_{\ell}v\left(\ell\right)$, where the sum is
over all $\left(m,n\right)$ framed lattices $\ell$. \label{thm: sum of vs} 
\label{section: counting polygon free mosaics}
\end{thm}

Throughout this section, let $v$ be as in Theorem~\ref{thm: sum of vs}.
Thus, for example, if $\ell$ is the framed lattice on the right of
Figure~\ref{fig:example of f mapping}, we have $v\left(\ell\right)=-7^{13}$,
as there are no $\begin{smallmatrix}0 & 1\\
1 & 0
\end{smallmatrix}$ or $\begin{smallmatrix}1 & 0\\
0 & 1
\end{smallmatrix}$ cells, twelve $\begin{smallmatrix}0 & 0\\
0 & 0
\end{smallmatrix}$ cells, one $\begin{smallmatrix}1 & 1\\
1 & 1
\end{smallmatrix}$ cell, two $\begin{smallmatrix}0 & 0\\
1 & 0
\end{smallmatrix}$ cells, and one $\begin{smallmatrix}1 & 0\\
1 & 1
\end{smallmatrix}$ cell.

The following result, which is an important part of proving Theorem~\ref{thm: sum of vs},
connects local and global information in the sense that $v\left(\ell\right)$
is calculated by multiplying values for each cell of $\ell$, whereas
the number of polygons is a characteristic of an entire mosaic.
\begin{prop}
For a polygon mosaic $P$, let $\ell=\mathscr{L}\left(P\right)$.
If $P$ contains an even {[}resp. odd{]} number of polygons then $v\left(\ell\right)$
is positive {[}resp. negative{]}, where $v$ is as in Theorem~\ref{thm: sum of vs}.
\label{prop: sign of v} 
\end{prop}

\begin{proof}
{[}{[}Insert proof of this proposition here.{]}{]}
\end{proof}
%
\begin{proof}[Proof of Theorem~\ref{thm: sum of vs}]
 Fixing $m$ and $n$, throughout this proof we use ``polygon mosaic''
for ``$\left(m,n\right)$ polygon mosaic,'' ``framed lattice''
for ``$\left(m,n\right)$ framed lattice,'' and so on.

We introduce the following identification for polygon mosaics. Let
$\left\{ P_{i}\right\} _{i\in J}$ be the set of all polygon mosaics
that have exactly one polygon, and for any subset $I$ of $J$, let
$P_{I}$, if it exists, be the polygon mosaic that contains exactly
the polygons of $P_{i}$ for all $i\in I$. For example, $P_{\emptyset}$
is the polygon mosaic containing no polygons, i.e., the mosaic entirely
composed of blank tiles, and for any $i\in J$, $P_{\left\{ i\right\} }$
always exists and equals $P_{i}$. For an arbitrary $I\subseteq J$,
$P_{I}$ exists if and only if the polygons of $P_{i}$, for $i\in I$,
do not ``overlap.'' More precisely, $P_{I}$ exists if and only
if for every mosaic location, there is at most one $i\in I$ such
that $P_{i}$ has a non-blank tile there. Notice that every polygon
mosaic can be uniquely written as $P_{I}$ for some $I\subseteq J$.

Next we look at some sets of mosaics. For any $I\subseteq J$ such
that $P_{I}$ exists, let $M_{I}$ be the set of mosaics that contain
the polygons of $P_{I}$, i.e., the set of mosaics that can be formed
from $P_{I}$ by using any tile wherever $P_{I}$ has a blank tile.
Notice that the size $\left|M_{I}\right|$ of this set equals $7^{a}$,
where $a$ is the number of blank tiles in $P_{I}$. By the definition
of $\mathscr{L}$, the number of cells of $\mathscr{L}\left(P_{I}\right)$
that are either $\begin{smallmatrix}0 & 0\\
0 & 0
\end{smallmatrix}$ or $\begin{smallmatrix}1 & 1\\
1 & 1
\end{smallmatrix}$ is also $a$, and since $\mathscr{L}\left(P_{I}\right)$ is proper,
i.e. it has neither $\begin{smallmatrix}0 & 1\\
1 & 0
\end{smallmatrix}$ nor $\begin{smallmatrix}1 & 0\\
0 & 1
\end{smallmatrix}$ cells, we see that $\left|v\left(\mathscr{L}\left(P_{I}\right)\right)\right|=\left|M_{I}\right|$.
Furthermore, $\left|I\right|$ is the number of polygons in $P_{I}$,
so Proposition~\ref{prop: sign of v} tells us that $v\left(\mathscr{L}\left(P_{I}\right)\right)$
and $\left(-1\right)^{\left|I\right|}$ have the same sign. Therefore,
if $P_{I}$ exists then 
\[
v\left(\mathscr{L}\left(P_{I}\right)\right)=\left(-1\right)^{\left|I\right|}\left|M_{I}\right|.
\]

We must show that the the sum $\sum_{\ell}v\left(\ell\right)$ over
all framed lattices $\ell$ is equal to number of polygon-free mosaics.
The sum is unchanged if we restrict to the sum over all proper framed
lattices $\ell$ because $v\left(\ell\right)=0$ if $\ell$ is not
proper. Furthermore, since $\mathscr{L}$ gives a bijection from polygon
mosaics to proper framed lattices, 
\[
\sum_{\ell}v\left(\ell\right)=\sum v\left(\mathscr{L}\left(P_{I}\right)\right)=\sum\left(-1\right)^{\left|I\right|}\left|M_{I}\right|,
\]
 where the second and third sums are over all $I\subseteq J$ such
that $P_{I}$ exists. The total number of mosaics, which happens to
be $7^{mn}$, equals $\left|M_{\emptyset}\right|$. Notice that if
$J=\emptyset$ if and only if $m=1$ or $n=1$. If this is the case
then every mosaic is polygon-free, the third sum above is $\left|M_{\emptyset}\right|$,
and our theorem is complete. Therefore we may assume that $m\ge2$
and $n\ge2$, so that $J\neq\emptyset$ and the third sum has a nonzero
term with $I\neq\emptyset$, for example, $I=\left\{ i\right\} $
for any $i\in J$.

For any $I\neq\emptyset$ such that $P_{I}$ exists, $M_{I}=\bigcap_{i\in I}M_{\left\{ i\right\} }$
by the definition of the $M_{I}$. Thus 
\[
\sum_{\ell}v\left(\ell\right)=\left|M_{\emptyset}\right|-\sum\left(-1\right)^{\left|I\right|-1}\left|\bigcap_{i\in I}M_{\left\{ i\right\} }\right|,
\]
 where the sum on the right is over all $I$ such that $P_{I}$ exists
and $\emptyset\neq I\subseteq J$. In this equation, the value on
the right side is unchanged if instead we sum over all $I$ such that
$\emptyset\neq I\subseteq J$. This is because when $P_{I}$ is not
defined, the $P_{i}$, $i\in I$, ``overlap'' in the sense described
above, and so $\bigcap_{i\in I}M_{\left\{ i\right\} }=\emptyset$.
By the inclusion-exclusion principle (see the Appendix), this summation
over all $I$ with $\emptyset\neq I\subseteq J$ equals $\left|\bigcup_{i\in J}M_{\left\{ i\right\} }\right|$,
i.e., the number of mosaics that contain at least one polygon. Since
$\left|M_{\emptyset}\right|$ is the total number of mosaics, we conclude,
as desired, that $\sum_{\ell}v\left(\ell\right)$ is the number of
polygon-free mosaics.
\end{proof}

\end{document}
